\scanpage{155}
Normalité des extensions quadratiques \quad $A[S]/(S^2-a)=B$

On suppose d'abord $A$ anneau jauge, soit $p$ sa car. résiduelle.
% slanted
\margin{\uline{NB} Si $p\ne2$, alors $B$ est local sss \uline{ou} $=a\in\mathfrak m$ \uline{ou} ($a\deleted{\in}\notin\mathfrak m$ et) \deleted{\unsure{$a\in$}} $\dot a\notin k^2$}
$p\ne2$, alors $B$ normal sss $a\notin\mathfrak m^2$ -- et cette condition est nécessaire en tous cas (y compris $p=2$). Si $p=2$, distinguons deux cas:

a) $a\in\mathfrak m$ alors $B$ normal sss $a\notin\mathfrak m^2$

b) $a\notin\mathfrak m$ i.e. $u$ une unité. Soit $\dot a$ son image dans $K'=k(A)=A/\mathfrak m$.

$1^\circ$) $\dot a\notin k^2$, alors $B$ est normal (c'est un anneau jauge de corps résiduel $k[\sqrt{\dot a}]$)

$2^\circ$) $\dot a\in k^2$, soit $\dot a=\dot\alpha^2$, où $\alpha\in A$ déterminé mod $\deleted{\unsure{\mathfrak m}}\unsure{\mathfrak m}$ -- donc $a=\alpha^2\deleted{\unsure{+}}-b$, avec $b\inserted{=\alpha^2-a}\in\mathfrak m$. Posons $T=S-\alpha$, alors
% slanted
\margin{\uline{NB} $\widetilde b\deleted{\unsure{\in}}\in\mathfrak m/\mathfrak m^2$ ne dépend que de $a$, et pas du choix de $\alpha$.}
\[
\begin{aligned}
S^2-a&=(T+\alpha)^2-a=T^2+2\alpha T+b\\
B&\simeq A[T]/(T^2+2\alpha T+b),\quad
\widehat B\simeq\underbrace{A[[T]]}_{C}/(T^2+2\alpha T+b)
\end{aligned}
\]
% slanted
\margin{NB $2\in\mathfrak m\subset\mathfrak r(C)$ \quad $T\in\mathfrak r(C)$}
$T^2,2\alpha T$ sont dans $\deleted{\unsure{\mathfrak m|M}}\mathfrak r(C)^2$, donc $T^2+2\alpha T+b\notin(\mathfrak r(C))^2$ sss $b\notin\mathfrak r(C)^2$ i.e $b\notin\mathfrak m^2$, i.e. $\widetilde b\ne0$ dans $\mathfrak m/\mathfrak m^2$

Donc les cas de normalité sont les suivants
\[
\begin{aligned}
&a\in\mathfrak m\setminus\mathfrak m^2,\ \text{\uline{ou}}\
a\ \text{de la forme }\alpha^2-b\ \text{avec }\alpha\ \text{unité},\
b\in\mathfrak m\setminus\mathfrak m^2\\
&\text{\uline{ou}}\quad a\notin\mathfrak m,\quad\dot a\notin k^2
\end{aligned}
\]
% slanted
\margin{Donc les cas où $B$ \uline{non} normal sont ceux où $a\in\mathfrak m^2$, \uline{ou} $\unsure{\widehat a|\dot a|a}\bmod\mathfrak m^2$ le carré d'une unité de $A/\mathfrak m^2$. \uline{Ce dernier cas \unsure{correspond|correspondra} à une singularité quadratique ordinaire.}}

\uline{Application} aux rev. quadratiques de $\mathbb P^1_{\mathbb Z}$ \inserted{(ou de $\mathbb E^1_{\mathbb Z}$)} génériquement \deleted{\unsure{non|normales}} étales sur chaque fibre \inserted{au} dessus de $\mathrm{Spec}\,\mathbb Z$.

Le revêtement \deleted{\unsure{au dessus du pt}}\inserted{quadratique} du pt gén. correspond à une ext. quadratique de $\mathbb Q(t)$, donc décrit par un élément de $\mathbb Q(t)^*/\mathbb Q(t)^{\unsure{*|2|*2}}$.

Or, $\mathbb Q(t)^{*}/\unsure{\{\pm1\}}$ a une base formée des éléments suivants

a) Les nbs premiers $p\in\mathbb P$ \unsure{!} $(\mathbb P=\{2,3,5,7\ldots\})$

b) Les polynomes $P=a_0t^n+a_1t^{n-1}+\ldots+a_n\in\mathbb Z[t]$, avec: $a_0>\deleted{\unsure{1}}0$, \deleted{\unsure{et}} les $a_i$ premiers entre eux (i.e. $\deleted{\unsure{\forall}}\forall p\in\mathbb P$, $P\bmod p\in\mathbb Z_p[t]$ \deleted{est} $\ne0$), $P$ premier dans $\mathbb Q[T]$.

Ainsi un élément de $\mathbb Q\unsure{[t]|(t)|(t)^*}/\mathbb Q(t)^{\unsure{*|2|*2}}$ est représenté de façon unique par un produit
\[
\unsure{P|\mathbb P}=\underbrace{\pm}_{\varepsilon}\,
\underbrace{p_1p_2\cdots p_{\unsure{h|k|n}}}_{q}\,
\underbrace{P_1\cdots P_\ell}_{Q}\tag*{(1)}
\]
où les $p_i\in\mathbb P$ sont distincts, et les $P_j$ comme dans b) aussi ! Cela signifie ainsi qu'on aura
\[
P=\pm qQ\qquad
\left\{\begin{aligned}
&q\in\mathbb N\ \text{``quadratfrei''}\\
&Q\in\mathbb Z[t]\ \text{à coeff. dominant }>0\\
&\text{``quadratfrei'' dans }\mathbb Q[t]\\
&\text{à coeff. premiers entre eux.}
\end{aligned}\right.
\]
% archive sheet number: 154
